% Erratum E7 for inclusion in main.tex.
% Suggested use: % Erratum E7 for inclusion in main.tex.
% Suggested use: % Erratum E7 for inclusion in main.tex.
% Suggested use: % Erratum E7 for inclusion in main.tex.
% Suggested use: \input{erratum_E7_gamma_positive}

\subsection*{Erratum E7: the restriction $\gamma>0$ when $\nu/\gamma$ appears}
\label{s:erratum:E7}

Several definitions and statements use the quotient $\nu/\gamma$.  They
should therefore be read with the restriction $\gamma>0$, unless a separate
convention is explicitly introduced for the degenerate case $\gamma=0$.

In particular, after the definition of $X_\gamma$ in equation (21), the
interpolation spaces should be introduced as follows:
\[
        \text{for } \gamma>0 \text{ and } 0\leq \nu\leq \gamma,
        \qquad
        X_{\nu:\gamma}:=(X,X_\gamma)_{\nu/\gamma,\infty}.
\]
The case $\gamma=0$ is not covered by this formula, since then necessarily
$\nu=0$ but the expression $\nu/\gamma$ is the indeterminate quotient $0/0$.
If desired, one may add the separate convention
\[
        X_{0:0}:=X,
        \qquad
        \norm{\cdot}{0:0}:=\norm{\cdot}{}.
\]
This convention is not needed for the subsequent arguments.

The same restriction is needed in Proposition 2.4.  Its statement
should begin
\[
        \text{Let } \gamma>0,\quad 0\leq \nu\leq \gamma,
        \quad\text{and }x\in X.
\]
Indeed, the displayed identity in Proposition 2.4 contains the three quantities
\[
        \norm{(T^*T)^\gamma\omega}{}^{-\nu/\gamma},
        \qquad
        N_{\nu/\gamma},
        \qquad
        \norm{x}{\nu:\gamma},
\]
all of which are undefined as written when $\gamma=0$.

With the correction $\gamma>0$, the proposition itself is unchanged.  The
endpoint cases discussed in its proof are then the meaningful endpoints
$\nu=0$ and $\nu=\gamma$ for a fixed positive $\gamma$.  If one separately
adopts the convention $X_{0:0}=X$, then the missing degenerate case
$\gamma=\nu=0$ reduces only to the elementary identity
\[
        \sup_{\norm{\omega}{}=1}|\scalar{x,\omega}{}|
        =
        \norm{x}{}.
\]
Thus this erratum is a domain-of-notation correction: it does not change any
of the nondegenerate estimates or their proofs.


\subsection*{Erratum E7: the restriction $\gamma>0$ when $\nu/\gamma$ appears}
\label{s:erratum:E7}

Several definitions and statements use the quotient $\nu/\gamma$.  They
should therefore be read with the restriction $\gamma>0$, unless a separate
convention is explicitly introduced for the degenerate case $\gamma=0$.

In particular, after the definition of $X_\gamma$ in equation (21), the
interpolation spaces should be introduced as follows:
\[
        \text{for } \gamma>0 \text{ and } 0\leq \nu\leq \gamma,
        \qquad
        X_{\nu:\gamma}:=(X,X_\gamma)_{\nu/\gamma,\infty}.
\]
The case $\gamma=0$ is not covered by this formula, since then necessarily
$\nu=0$ but the expression $\nu/\gamma$ is the indeterminate quotient $0/0$.
If desired, one may add the separate convention
\[
        X_{0:0}:=X,
        \qquad
        \norm{\cdot}{0:0}:=\norm{\cdot}{}.
\]
This convention is not needed for the subsequent arguments.

The same restriction is needed in Proposition 2.4.  Its statement
should begin
\[
        \text{Let } \gamma>0,\quad 0\leq \nu\leq \gamma,
        \quad\text{and }x\in X.
\]
Indeed, the displayed identity in Proposition 2.4 contains the three quantities
\[
        \norm{(T^*T)^\gamma\omega}{}^{-\nu/\gamma},
        \qquad
        N_{\nu/\gamma},
        \qquad
        \norm{x}{\nu:\gamma},
\]
all of which are undefined as written when $\gamma=0$.

With the correction $\gamma>0$, the proposition itself is unchanged.  The
endpoint cases discussed in its proof are then the meaningful endpoints
$\nu=0$ and $\nu=\gamma$ for a fixed positive $\gamma$.  If one separately
adopts the convention $X_{0:0}=X$, then the missing degenerate case
$\gamma=\nu=0$ reduces only to the elementary identity
\[
        \sup_{\norm{\omega}{}=1}|\scalar{x,\omega}{}|
        =
        \norm{x}{}.
\]
Thus this erratum is a domain-of-notation correction: it does not change any
of the nondegenerate estimates or their proofs.


\subsection*{Erratum E7: the restriction $\gamma>0$ when $\nu/\gamma$ appears}
\label{s:erratum:E7}

Several definitions and statements use the quotient $\nu/\gamma$.  They
should therefore be read with the restriction $\gamma>0$, unless a separate
convention is explicitly introduced for the degenerate case $\gamma=0$.

In particular, after the definition of $X_\gamma$ in equation (21), the
interpolation spaces should be introduced as follows:
\[
        \text{for } \gamma>0 \text{ and } 0\leq \nu\leq \gamma,
        \qquad
        X_{\nu:\gamma}:=(X,X_\gamma)_{\nu/\gamma,\infty}.
\]
The case $\gamma=0$ is not covered by this formula, since then necessarily
$\nu=0$ but the expression $\nu/\gamma$ is the indeterminate quotient $0/0$.
If desired, one may add the separate convention
\[
        X_{0:0}:=X,
        \qquad
        \norm{\cdot}{0:0}:=\norm{\cdot}{}.
\]
This convention is not needed for the subsequent arguments.

The same restriction is needed in Proposition 2.4.  Its statement
should begin
\[
        \text{Let } \gamma>0,\quad 0\leq \nu\leq \gamma,
        \quad\text{and }x\in X.
\]
Indeed, the displayed identity in Proposition 2.4 contains the three quantities
\[
        \norm{(T^*T)^\gamma\omega}{}^{-\nu/\gamma},
        \qquad
        N_{\nu/\gamma},
        \qquad
        \norm{x}{\nu:\gamma},
\]
all of which are undefined as written when $\gamma=0$.

With the correction $\gamma>0$, the proposition itself is unchanged.  The
endpoint cases discussed in its proof are then the meaningful endpoints
$\nu=0$ and $\nu=\gamma$ for a fixed positive $\gamma$.  If one separately
adopts the convention $X_{0:0}=X$, then the missing degenerate case
$\gamma=\nu=0$ reduces only to the elementary identity
\[
        \sup_{\norm{\omega}{}=1}|\scalar{x,\omega}{}|
        =
        \norm{x}{}.
\]
Thus this erratum is a domain-of-notation correction: it does not change any
of the nondegenerate estimates or their proofs.


\subsection*{Erratum E7: the restriction $\gamma>0$ when $\nu/\gamma$ appears}
\label{s:erratum:E7}

Several definitions and statements use the quotient $\nu/\gamma$.  They
should therefore be read with the restriction $\gamma>0$, unless a separate
convention is explicitly introduced for the degenerate case $\gamma=0$.

In particular, after the definition of $X_\gamma$ in equation (21), the
interpolation spaces should be introduced as follows:
\[
        \text{for } \gamma>0 \text{ and } 0\leq \nu\leq \gamma,
        \qquad
        X_{\nu:\gamma}:=(X,X_\gamma)_{\nu/\gamma,\infty}.
\]
The case $\gamma=0$ is not covered by this formula, since then necessarily
$\nu=0$ but the expression $\nu/\gamma$ is the indeterminate quotient $0/0$.
If desired, one may add the separate convention
\[
        X_{0:0}:=X,
        \qquad
        \norm{\cdot}{0:0}:=\norm{\cdot}{}.
\]
This convention is not needed for the subsequent arguments.

The same restriction is needed in Proposition 2.4.  Its statement
should begin
\[
        \text{Let } \gamma>0,\quad 0\leq \nu\leq \gamma,
        \quad\text{and }x\in X.
\]
Indeed, the displayed identity in Proposition 2.4 contains the three quantities
\[
        \norm{(T^*T)^\gamma\omega}{}^{-\nu/\gamma},
        \qquad
        N_{\nu/\gamma},
        \qquad
        \norm{x}{\nu:\gamma},
\]
all of which are undefined as written when $\gamma=0$.

With the correction $\gamma>0$, the proposition itself is unchanged.  The
endpoint cases discussed in its proof are then the meaningful endpoints
$\nu=0$ and $\nu=\gamma$ for a fixed positive $\gamma$.  If one separately
adopts the convention $X_{0:0}=X$, then the missing degenerate case
$\gamma=\nu=0$ reduces only to the elementary identity
\[
        \sup_{\norm{\omega}{}=1}|\scalar{x,\omega}{}|
        =
        \norm{x}{}.
\]
Thus this erratum is a domain-of-notation correction: it does not change any
of the nondegenerate estimates or their proofs.
